\chapter{Discussion} \label{ch:discussion}

\section{Impact of Constrained Decoding vs.\ Model Scale}

The 71.25 percentage point gap between Config~B and Config~CD is almost
entirely a format compliance gap. Of the 394 failures in Config~B, virtually
all are structural: the model produces conversational text, multiple JSON
objects, or empty output rather than a single valid function call. The six
correct outputs in Config~B are semantically accurate but happen to be
formatted correctly by chance. The model is not unable to reason about the
queries; it cannot express its answers in the required structured format
without guidance.

This distinction has a direct consequence for interpreting the frontier
comparison. Config~CD achieves 72.75\% on BFCL~v4 \texttt{simple\_python},
matching GPT-4.1 (72.67\%) and Claude Sonnet~4.5 (72.58\%), and exceeding
Claude Haiku~4.5 (71.00\%)~\cite{bfcl2024}. The primary mechanism is that
constrained decoding enforces valid JSON and valid function names by
construction, which is precisely what frontier models' native tool-calling
APIs also do. Once format compliance is equalized, a 7B model and a frontier
model are distinguished only by argument value accuracy, where a smaller
model can remain competitive on straightforward single-function queries.

The comparison is specific to the \texttt{simple\_python} category. Each test
case provides exactly one candidate function, eliminating function selection
as a source of error. With one candidate and \texttt{guided\_choice},
selection accuracy is 100\% regardless of model size. The remaining 27.25\%
failure rate is entirely in argument values: wrong optional parameter defaults,
numeric precision errors, and string format mismatches. Frontier models make
the same class of errors; the gap of 6.83 percentage points to Gemini~3 Pro
(79.58\%) reflects a higher frequency of correct value conventions in larger
models, not a qualitatively different failure mode.

The implication for RQ1 is that the raw unoptimized gap of approximately
98.5 percentage points (Config~B vs.\ frontier) collapses to approximately
7 percentage points with constrained decoding alone, on this benchmark
category. Scale is not the bottleneck for format compliance. The residual
gap is a semantic problem, attributable to the model's learned associations
between function schema contexts and specific value forms.

\section{Cost-Benefit Analysis of Each Method}

Constrained decoding yields the highest return of any technique evaluated.
A 71.25 percentage point accuracy gain is achieved with no additional
inference cost beyond the guided decoding schema: no extra parameters, no
additional tokens, and no training. The gain is entirely structural: logit
masking at each decoding step eliminates the malformed output that accounts
for 98.5\% of failures in Config~B~\cite{willard2023efficient, kwon2023vllm}.

AWQ INT4 quantization adds a further infrastructure benefit at near-zero
accuracy cost~\cite{lin2024awq}. The $-0.5$ percentage point change (two
test cases) falls within run-to-run variance. Memory is reduced by 63.5\%
(14.25~GiB to 5.20~GiB) and model loading time by 83.5\% (212~s to 35~s).
The cost-benefit ratio is strongly positive: substantial deployment benefits
at negligible accuracy cost.

Few-shot prompting (Config~PE) and chain-of-thought reasoning
(Config~CD+Q+ITC) both produce negative outcomes~\cite{brown2020language,
wei2022chain}. Config~PE reduces accuracy by 2.5 percentage points with a
context overhead of approximately 300 tokens per query. Config~CD+Q+ITC
reduces accuracy by 6.5 percentage points while approximately doubling
per-case latency. Both techniques add cost and reduce accuracy. The failure
mechanism is the same in both cases: the residual 27\% failure rate after
constrained decoding is caused by a mismatch between the model's learned
value priors and the benchmark's ground truth conventions. Prompting
interventions that ask the model to reason more explicitly or reference
examples move it further from the benchmark's arbitrary conventions rather
than toward them. Two independent negative results with the same directional
outcome and a shared mechanistic explanation establish that prompting cannot
address this failure class.

Retrieval-augmented generation produces the largest accuracy drop: $-24$
percentage points relative to Config~CD+Q, despite a retrieval recall@5 of
97.2\%~\cite{lewis2020retrieval}. The bottleneck is not retrieval quality
but candidate disambiguation. When multiple semantically similar functions
are retrieved, the 7B model cannot reliably identify the one matching the
user's intent. The infrastructure cost of RAG is substantial (an offline
FAISS index and an embedder at query time), making its negative accuracy
outcome particularly unfavorable.

LoRA fine-tuning on the \texttt{xlam-function-calling-60k}
dataset~\cite{zhang2024xlam, hu2022lora} produces a regression of 3
percentage points under Config~CD+FT relative to Config~CD. The regression
is attributed to a format mismatch between the training distribution, which
encodes assistant turns as Python call syntax, and the evaluation pipeline,
which expects a JSON argument object. The technique is not ruled out by this
result: the aligned v2 ablation (Config~CD+FT-aligned) is pending. Config~FT-only
confirms that fine-tuning does improve format compliance in the unguided
setting ($+12.25$ percentage points over Config~B), establishing that the
training signal is meaningful. The 56 percentage point gap between
Config~FT-only and Config~CD+FT confirms that constrained decoding remains
the dominant contributor even after fine-tuning.

Taken together, the results support a clear ordering by cost-benefit ratio
for this task: constrained decoding first, quantization second, and
training-based methods as the remaining path to accuracy above the
no-training ceiling. Prompting interventions offer no benefit on this
failure class and add latency or context overhead.

\section{Limitations}

The results are based on the \texttt{simple\_python} category of BFCL~v4,
which is the least demanding of the four BFCL categories~\cite{bfcl2024}.
Each test case consists of a single user turn, a single candidate function,
and a single expected call. The constrained decoding advantage depends on
the oracle function provision: \texttt{guided\_choice} guarantees correct
function selection only when the correct function is the sole candidate.
The RAG results (47.75\% at top-5) show that this assumption fails under
realistic conditions. Generalisation to harder BFCL categories (multiple
function candidates, parallel function calls, multi-turn interactions) is
not established and may not hold.

The full configuration ladder is evaluated only on the 7B model. Smaller
Qwen~2.5 sizes (0.5B, 1.5B, 3B) were not evaluated across all
configurations, so the interaction between model scale and each optimization
technique is not characterised. Whether constrained decoding closes the
frontier gap equally at smaller scales, or whether there is a scale threshold
below which semantic failures dominate, remains an open question.

The regression in Config~CD+FT is attributed to the format mismatch between
the xlam training distribution and the evaluation pipeline. This attribution
is supported by failure pattern analysis and the v1 training format design,
but has not yet been confirmed by a fully controlled experiment. The aligned
ablation (Config~CD+FT-aligned), which holds all other variables fixed and
changes only the training output format, will directly test this hypothesis.
Until those results are available, the attribution remains an explanation
rather than a confirmed finding.

Multi-step agentic performance is not yet reported. The $\tau$-bench
evaluation, which requires the model to reason across multiple tool calls
with real observations, is ongoing. Simple function calling accuracy on
BFCL does not predict multi-step performance; the 7 percentage point gap
between Config~CD and Gemini~3 Pro on the simple category may widen
substantially on multi-turn tasks.

The BFCL ground truth encodes specific value conventions that are often
arbitrary from a deployment perspective. A model that produces
\texttt{"gasoline"} where the ground truth expects \texttt{"gas"} is scored
as incorrect despite both being valid in a real API call. The 27\% residual
failure rate after constrained decoding may overestimate the practical error
rate for production use, where value conventions are defined by the API
schema rather than by a benchmarking dataset.

\section{Practical Implications}

Config~CD+Q addresses RQ3 directly. A 72.25\% AST accuracy is achieved on
5.20~GiB of GPU memory with a 35~s model load time on an NVIDIA L40S. The
same configuration fits on a consumer RTX~4090 (24~GiB) with ample remaining
capacity for the key-value cache. At near-zero accuracy cost relative to full
precision, Config~CD+Q is the production-viable configuration for the
\texttt{simple\_python} use case on consumer hardware.

For a cascade architecture where an SLM handles routine tool-calling
requests and escalates difficult ones to a frontier model, Config~CD+Q
defines the SLM operating point. Approximately 73\% of requests are handled
correctly at 5.20~GiB memory; the remaining 27\% are escalated. The
escalation boundary corresponds to the residual failure classes: wrong
optional parameter defaults, numeric precision mismatches, and string format
conventions. Queries involving domain-specific numeric precision or
underspecified optional parameters are candidates for escalation; queries
with explicit argument values and unambiguous types are not.

Chain-of-thought reasoning is contraindicated at the SLM tier. Config~CD+Q+ITC
raises the failure rate by 6.5 percentage points relative to Config~CD+Q,
increasing escalation traffic rather than reducing it~\cite{wei2022chain}.
The mechanism is structural: the reasoning step moves the model toward
linguistically natural value forms that differ from the conventions the
ground truth encodes.

For tool catalogs with many functions, RAG-based tool selection requires a
disambiguation capability that the 7B model does not have at top-5
retrieval~\cite{lewis2020retrieval}. A more viable configuration would use
top-1 or top-2 retrieval, a re-ranking step before the selection stage, or
a fine-tuning signal on tool-selection examples to reduce the frequency of
sibling-function confusion.

Fine-tuning remains the most direct path to accuracy above the 72--73\%
no-training ceiling. The key constraint is training data alignment: the
training format must match the inference pipeline's prompt structure and
output conventions exactly. The v1 xlam experiment shows that a
general-purpose function-calling dataset with a different output convention
actively degrades performance under constrained decoding. A production
fine-tuning pipeline would require training examples generated from the
target deployment's own tool schemas and argument conventions, not from a
general-purpose corpus.
